% Junior Level 1, Week 8 -- answer key
\answersweek{8}{Draw With the Pen}{Answers for workbook chapter 8 and play-at-home card 8.}

\section{Lesson questions}
\begin{answer}[Exit question]
If the turn is taken out of the repeat, the cat draws one long straight line (100 steps, four times, until it stops at the edge of the stage) and then turns once at the end. It is no longer a square.
\end{answer}

\begin{answer}[Reflection]
Any drawing of a square, a ring of cats or a free drawing is right. \enquote{Pen down makes the cat draw} (or \enquote{the cat leaves a line}, \enquote{the crayon goes on the paper}) shows understanding.
\end{answer}

\begin{answer}[Password question (safety talk)]
Say no, and tell a grown-up. A password is like a toothbrush: you never share it, not even with a best friend. Only the grown-ups who look after you may know it.
\end{answer}

\section{Workbook exercises}
\begin{answer}[Exercise 8.1: Draw a Square]
\textbf{What you should see:} after the green flag, the cat gets small, old lines disappear, and the cat draws a square: it walks, makes a quarter turn to the right, four times, and ends where it started. Clicking the flag again draws the same square in the same place.

\textbf{What to look for:} \tickbox{} the Pen extension is added \quad \tickbox{} \enquote{pen down} is above the repeat \quad \tickbox{} move 100 and turn 90 are both inside \enquote{repeat 4} \quad \tickbox{} the child counts four sides.
\end{answer}

\begin{answer}[Exercise 8.2: The One-Side Square]
The circled block is \enquote{turn 90 degrees}: it is under the repeat, not inside it. So the cat moves 100 steps four times in a row (one long line that stops at the edge) and turns only once, at the end. Fix: drag the turn up into the repeat, under the move. Then the cat draws all four sides.
\end{answer}

\begin{answer}[Exercise 8.3: Stamp a Ring]
\textbf{What you should see:} six cats in a ring, each one turned a little more than the one before, like the petals of a flower. The real cat sits on top of the first stamp. The turn 60 comes from the turn card (hexagon, 6 sides).

\textbf{What to look for:} \tickbox{} stamp, move and turn are all inside \enquote{repeat 6} \quad \tickbox{} \enquote{erase all} at the top clears old stamps.
\end{answer}

\section{Play-at-home card}
\begin{answer}[My greeting card idea]
Any person or pet and any drawing is right. Good ideas name who the card is for and one thing that happens (\enquote{a cat dances for Grandma}, \enquote{balloons for my friend's birthday}). Keep the drawing: it becomes box 1 of the storyboard in week 9. Use a nickname or \enquote{Grandma}, \enquote{my friend}, never a full real name.
\end{answer}

\begin{answer}[Walk a square]
After four times \enquote{walk 3 steps, quarter turn} the child is back where they started, facing the same way. On paper, four sides with a quarter turn at each corner make a square.
\end{answer}
